\section{Conclusion}
% 结论。
\label{sec:conclusion}
We presented \sys, a policy runtime that models each GPU resource as an asynchronous state machine whose state transitions are directed by verified eBPF programs across the host-device boundary.
% 我们提出了 \sys,一个策略运行时,它将每个 GPU 资源建模为异步状态机,其状态迁移由跨越主机-设备边界的经验证的 eBPF 程序驱动。
Across four workloads, policies implemented with \sys{} improve throughput by up to $1.76\times$ over application-tuned baselines and reduce tail latency by up to $2\times$, demonstrating that OS kernel extensibility extends effectively to GPU management.
% 在四类工作负载上,通过 \sys{} 实现的策略相比应用调优的基线将吞吐量最高提升 $1.76\times$、将尾延迟最高降低 $2\times$,表明 OS 内核可扩展性可以有效延伸到 GPU 管理。
For seven published resource-management policies, the metrics in \Cref{fig:matched-ports} differ by less than 2\% between \sys{} and native implementations.
% 对于七种已发表的资源管理策略，\Cref{fig:matched-ports} 中各项指标在 \sys{} 实现与原生实现之间的相对差异均小于 2\%。
